\par\medskip
\textls[100]{1. Nullstellen von \(U\) und \(T\) ausserhalb \((R)\).} Es sei
\[
f(z)\equiv z^n+Az^{n-1}+\cdots+C=T(r,\varphi)+iU(r,\varphi)=0,\quad z=x+iy,
\]
die vorgegebene algebraische Gleichung, wo \(r,\varphi\) Polarkoordinaten, \(x,y\) kartesische
Koordinaten sind. Um zu beweisen, dass die Kurven \(T=0,\ U=0\)
sich schneiden, ist es nötig, den Verlauf dieser Kurven zu untersuchen. Wir
benutzen dabei zuerst Polarkoordinaten.
\par\clearpage\SourceLeaf{319}\ArithmeticHead{10}{L}{A. OSTROWSKI, ÜBER DEN ERSTEN UND VIERTEN GAUSSSCHEN BEWEIS USW.}
Es sei \(\omega=\frac{\pi}{4n}\). Teilen wir durch \(4n\) Halbstrahlen \(\varphi=-\omega,\ \varphi=\omega\),
\(\varphi=3\omega\) usw. die ganze Ebene in \(4n\) Winkelräume und bezeichnen sie in der
entsprechenden Reihenfolge durch \((1),(2),\ldots,(4n)\), so gibt es, wie wir jetzt
beweisen werden, eine Zahl \(R\) mit den folgenden Eigenschaften:
\par
1) \(f(z)\) hat keine Wurzel vom absoluten Betrage \(\geqq R\).
\par
2) Man beschreibe um den Nullpunkt eine Kreislinie \((\overline R)\) mit dem Halbmesser
\(\overline R\geqq R\) und betrachte \(T\) und \(U\) auf dieser Kreislinie. Dann sind in den \(\varphi\)-Intervallen
\((4n+1)\) \(T\) und \(\frac{\partial U}{\partial\varphi}\) positiv, \(U\) ist hingegen zuerst \(<0\), dann \(>0\), hat also eine
Wurzel innerhalb jedes \(\varphi\)-Intervalls \((4n+1)\) auf jedem Kreise \((\overline R)\) und zwar nur eine,
wegen \(\frac{\partial U}{\partial\varphi}>0\). In den \(\varphi\)-Intervallen \((4n+2)\) ist \(U>0,\ \frac{\partial T}{\partial\varphi}<0\), \(T\) zuerst
\(>0\), dann \(<0\), also innerhalb jedes solchen Intervalls \((4n+2)\) einmal und nur einmal
\(=0\). In den \(\varphi\)-Intervallen \((4n+3)\) ist \(T<0,\ \frac{\partial U}{\partial\varphi}<0\), \(U\) wechselt je einmal
das Vorzeichen. In den \(\varphi\)-Intervallen \((4n)\) endlich ist \(U<0,\ \frac{\partial T}{\partial\varphi}>0\),
\(T\) wechselt je einmal das Vorzeichen.
\par
Um ein solches \(R\) zu finden, bezeichnen wir die absoluten Beträge der
Koeffizienten \(A,\ldots,C\) von \(f(z)\) mit \(a,\ldots,c\) und die Summe dieser absoluten
Beträge mit \(S\). Dann besitzt \(R=\sqrt{2}S+1\) die geforderte Eigenschaft. In
der Tat gilt mit diesem Wert von \(R\) für \(\overline R\geqq R\):
\[
\begin{aligned}
|T-\overline R^n\cos n\varphi|&\leqq a\overline R^{n-1}+\cdots+c\leqq\overline R^{n-1}S<\sqrt{\tfrac12}(\overline R)^n,\\
\left|\frac{\partial U}{\partial\varphi}-n\overline R^n\cos n\varphi\right|&\leqq(n-1)a\overline R^{n-1}+\cdots<n\overline R^{n-1}S<n\sqrt{\tfrac12}(\overline R)^n.
\end{aligned}
\]
\noindent Daher haben dann für jedes \(\varphi\), für das \(|\cos n\varphi|\geqq\sqrt{\tfrac12}\) ist, \(T\) und \(\frac{\partial U}{\partial\varphi}\) das Vorzeichen
von \(\cos n\varphi\). Dies findet also für alle \(\varphi\) statt, die für ganze \(k\) zwischen
\((8k-1)\omega\) und \((8k+1)\omega\), oder zwischen \((8k+3)\omega\) und \((8k+5)\omega\) liegen, d. h.
in den Intervallen \((1),(3),(5)\ldots\). Und zwar sind in \((1),(5),\ldots\) \(T\) und \(\frac{\partial U}{\partial\varphi}\) positiv,
in \((3),(7),\ldots\) negativ. Ganz analog sehen wir, dass, sobald \(|\sin n\varphi|\geqq\sqrt{\tfrac12}\) ist,
\(U\) für \(\overline R\geqq R\) das Vorzeichen von \(\sin n\varphi\) und \(\frac{\partial T}{\partial\varphi}\) das Vorzeichen von \(-\sin n\varphi\)
hat. Dies findet also für alle \(\varphi\) statt, die für ganze \(k\) zwischen \((8k+1)\omega\)
und \((8k+3)\omega\) oder zwischen \((8k+5)\omega\) und \((8k+7)\omega\) liegen, d. h. in \((2),(4)\),
\((6),\ldots\). Und zwar ist in \((2),(6),\ldots\) \(U\) positiv, \(\frac{\partial T}{\partial\varphi}\) negativ, in \((4),(8),\ldots\)
aber \(U\) negativ, \(\frac{\partial T}{\partial\varphi}\) positiv. — Daraus ergibt sich 1) und 2) sofort.
\par\clearpage\SourceLeaf{320}\ArithmeticHead{11}{R}{A. OSTROWSKI, ÜBER DEN ERSTEN UND VIERTEN GAUSSSCHEN BEWEIS USW.}
\HistoricalFoot{}{2\textsuperscript{*}}
\textls[100]{2. Verlauf von \(U=0\) und \(T=0\) ausserhalb \((R)\).} Danach ist es
leicht, den Verlauf etwa der Kurve \(U=0\) für \(r\geqq R\) zu beschreiben. Sie
zerfällt ausserhalb der Kreislinie \((R)\) mit dem Halbmesser \(R\) in \(2n\) Zweige, die
in den Intervallen \((1),(3),\ldots\) getrennt verlaufen. Fassen wir etwa die
Nullstellen von \(U\) ausserhalb \((R)\) in (1) ins Auge, so liegt auf jedem Kreise
\(r=\overline R(\overline R\geqq R)\) genau eine. Alle diese Nullstellen bilden aber eine stetige
Kurve. Um dies zu beweisen, brauchen wir nur Folgendes zu zeigen:
\par
Ist \(w\geqq R\) und \((w,\vartheta)\) die Nullstelle von \(U\), die in (1) auf dem Kreise
\(r=w\) liegt, so kann man jedem \(\varepsilon\) ein solches \(\delta\) zuordnen, dass, wenn \(w'\geqq R\),
\(w-\delta<w'<w+\delta\) ist und die auf dem Kreise \(r=w'\) in (1) liegende Wurzel
mit \((w',\vartheta(w'))\) bezeichnet wird, \(|\vartheta-\vartheta(w')|<\varepsilon\) bleibt.
\par
Da aber \(U\) in \((w,\vartheta)\) auf dem Kreise \(r=w\) von negativen zu positiven
Werten übergeht, so ist für ein hinreichend kleines \(\varepsilon'<\varepsilon\)
\[
U(w,\vartheta-\varepsilon')<0,\quad U(w,\vartheta+\varepsilon')>0,
\]
und sowohl \((w,\vartheta-\varepsilon')\) als auch \((w,\vartheta+\varepsilon')\) bleiben innerhalb (1). Da \(U\) stetig
in den Variabeln \(r,\varphi\) ist, so hat man für ein hinreichend kleines \(\delta\)
\[
U(w',\vartheta-\varepsilon')<0,\quad U(w',\vartheta+\varepsilon')>0,
\]
solange \(w-\delta<w'<w+\delta\) bleibt. Daher befindet sich auf dem Kreise \(r=w'\)
in (1) eine Nullstelle von \(U\), deren Argument zwischen \(\vartheta-\varepsilon'\) und \(\vartheta+\varepsilon'\) liegt,
für die also \(|\vartheta(w')-\vartheta|<\varepsilon\) ist.
\par
Und genau analog schliesst man für die übrigen Intervalle \((3),(5),\ldots\),
sowie für die Kurve \(T=0\).
\par
